% Chapter: Executive Summary
\chapter{Resumen Ejecutivo}

\begin{marketdatabox}[Snapshot del Mercado Minero Global]
\begin{itemize}[leftmargin=*]
    \item \textbf{Tamaño del Mercado Minero Global (2026):} \marketsize{2.16 trillones}
    \item \textbf{Proyección (2031):} \marketsize{2.93 trillones}
    \item \textbf{CAGR (2026--2031):} \growthrate{6.3}
    \item \textbf{Mercado Consultoría Minera (2025):} \marketsize{12.8 mil millones}
    \item \textbf{Proyección Consultoría (2034):} \marketsize{22.6 mil millones} | CAGR \growthrate{6.5}
    \item \textbf{Pipeline CAPEX Chile:} USD 70--105B (2026--2034) | 81\% brownfield
    \item \textbf{Pipeline CAPEX Perú:} USD 55B | inversión 2026: USD 6.3B
    \item \textbf{Pipeline CAPEX Brasil:} USD 45B | hierro + litio + tierras raras
\end{itemize}
\end{marketdatabox}

\section{Tesis de Inversión}

\begin{keyinsightbox}[Drivers Clave de Oportunidad]
\begin{enumerate}
    \item \textbf{Transición energética:} La demanda de cobre para electrificación, EVs e infraestructura de IA genera un déficit estructural de oferta proyectado para 2028--2030. Cada EV requiere $\sim$80 kg de cobre vs.\ $\sim$20 kg de un vehículo convencional.
    \item \textbf{Superciclo de inversión LATAM:} Chile, Perú y Brasil concentran más de USD 170B en pipeline minero. El CAPEX regional crece 8.4\% en 2026.
    \item \textbf{Minerales críticos y seguridad nacional:} EE.UU.\ ha movilizado >USD 30B en apoyo a producción doméstica de litio, cobre y tierras raras. Rusia invierte récord de USD 2.2B (Polyus, 2025).
    \item \textbf{Digitalización y ESG:} El mercado de Smart Mining crece de USD 17.2B (2026) a USD 31.4B (2034) — CAGR 7.81\%. Demanda creciente de consultoría especializada.
\end{enumerate}
\end{keyinsightbox}

\section{Hallazgos Principales}

\paragraph{Dinámica del Mercado.} La industria minera global alcanza USD 2.16T en 2026 \cite{researchandmarkets2026}, impulsada por la demanda de minerales de transición energética (Cu, Li, Ni, Co) \cite{iea2026}. Asia-Pacífico lidera en consumo; América Latina lidera en pipeline de inversión cuprífero \cite{bnamericas2026}.

\paragraph{Paisaje Competitivo.} El mercado de consultoría de ingeniería minera es fragmentado \cite{dataintelo2025}: los 5 principales actores (Worley \cite{worley2025}, Hatch, Jacobs, SRK, Ausenco) concentran $<$35\% del mercado. Las firmas boutique especializadas capturan nichos de alto valor en estudios de factibilidad, planificación minera y optimización operacional.

\paragraph{Riesgos Críticos.} El ``Efecto Seligdar'' (contratos rusos \cite{seligdar2025} que inflan el POM sin conversión a caja) ilustra el riesgo geopolítico sistémico. Las sanciones occidentales a Rusia han forzado una reorientación hacia China e India, creando complejidad operacional para consultoras internacionales.

\section{Recomendaciones Estratégicas}

\begin{recommendationbox}[Top 5 Recomendaciones Estratégicas]
\begin{enumerate}
    \item \textbf{Consolidar posición en Chile como mercado core:} Capitalizar el pipeline de USD 70--105B en cobre, con foco en brownfield (81\% del pipeline).
    \item \textbf{Expandir a Perú como segundo mercado prioritario:} Pipeline de USD 55B y USD 6.3B en inversión inmediata (2026).
    \item \textbf{Desarrollar capacidades digitales y ESG:} Posicionarse en el segmento de Smart Mining (CAGR 7.81\%) y consultoría de sostenibilidad.
    \item \textbf{Gestionar activamente el riesgo Rusia:} Aplicar el ``toggle Efecto Seligdar'' — separar POM de caja real; mantener exposición controlada vía intermediarios asiáticos.
    \item \textbf{Evaluar entrada en Brasil vía alianza estratégica:} Acceder al mercado de hierro/litio/tierras raras sin inversión masiva en infraestructura propia.
\end{enumerate}
\end{recommendationbox}
